\chnunumer{10536}
\chnuname{长沙理工大学}
\cclassnumber{TP391}
\cnumber{23408030281}
% \cnumber{*}                     % -------------- *
\csecret{公开}
\cmajor{工程硕士}  % 学位类别
\cdegreethesis{专业硕士学位论文}  % 自己的学位论文级别
\cdegreethesisfeiquan{非全日制专业硕士学位论文}
\cheading{硕士学位论文}  % 设置正文的页眉
\dtitle{基于 MLIR 的 Transformer 模型关键算子自动代码生成的研究}  %封面用论文标题，自己可手动断行\\
\ctitle{基于 MLIR 的 Transformer 模型关键算子自动代码生成的研究}  %页眉标题无需断行
\etitle{Research on Automatic Code Generation for Key Operators of Transformer Models Based on MLIR}
\caffil{计算机学院}  %学院名称
\csubjecttitle{学科专业}
\csubject{电子信息}  %学位领域
\cauthortitle{研究生}  % 学位
\cauthor{谢宝康}  %学生姓名
% \cauthor{*}  %学生姓名                     % -------------- *
\ename{XIE~Baokang}
% \ename{*}                     % -------------- *
\cbe{B.E.~(Hunan University of Humanities, Science and Technology)~2022}
\cdegree{thesis}
\cclass{Master of Engineering}
\ehnu{Changsha University of Science \& Technology}
\esupervisor{SONG Yun}
% \esupervisor{*}                     % -------------- *
\csupervisortitle{指导教师}
\csupervisor{宋云~~教授}  %导师姓名
% \csupervisor{*~~教授}  %导师姓名                     % -------------- *
\elevel{Professor}  %导师职称
\ocsupervisor{肖俊敏~~副研究员}  %校外导师
% \ocsupervisor{}                     % -------------- *
\cchair{李文军}
\ddate{2026年5月}  %论文答辩日期
\edate{April,~2026}

\untitle{长沙理工大学}
\declaretitle{学位论文原创性声明}
\declarecontent{
    本人郑重声明：所呈交的论文是本人在导师的指导下独立进行研究所取得的研究成果。除了文中特别加以标注引用的内容外，本论文不包含任何其他个人或集体已经发表或撰写的成果作品。对本文的研究做出重要贡献的个人和集体，均已在文中以明确方式标明。本人完全意识到本声明的法律后果由本人承担。
}
\authorizationtitle{学位论文版权使用授权书}
\authorizationcontent{
    本学位论文作者完全了解学校有关保留、使用学位论文的规定，同意学校保留并向国家有关部门或机构送交论文的复印件和电子版，允许论文被查阅和借阅。本人授权长沙理工大学可以将本学位论文的全部或部分内容编入有关数据库进行检索\scalebox{0.9}[1]{，}可以采用影印、缩印或扫描等复制手段保存和汇编本学位论文。同时授权中国科学技术信息研究所将本论文收录到《中国学位论文全文数据库》，并通过网络向社会公众提供信息服务。
}
\authorizationadd{本学位论文属于}
\authorsigncap{作者签名:}
\supervisorsigncap{导师签名:}
\signdatecap{日期:}


%\cdate{\CJKdigits{\the\year} 年\CJKnumber{\the\month} 月 \CJKnumber{\the\day} 日}
% 如需改成二零一二年四月二十五日的格式，可以直接输入，即如下所示
\cdate{2026年4月} %论文提交日期
% \cdate{\the\year 年\the\month 月 \the\day 日} % 此日期显示格式为阿拉伯数字 如2012年4月25日
%中文摘要%
\cabstract{
Transformer 架构在自然语言处理、代码生成及大模型推理等任务中得到了广泛使用，矩阵乘法（Matrix Multiplication, Matmul）和多头注意力机制（Multi-Head Attention，MHA）作为其关键算子（Kernel），它们的执行效率已经成为影响模型部署性能的主要因素。异构硬件平台在模型执行、内存层次结构、软件生态系统上的差异比较大，使用高性能库或手动编写优化算子的传统方法，会面临适应性复杂、开发成本高、性能可移植性不够等问题。同时，现在的深度学习编译器在块级（Block）逻辑表达、低级内存访问控制和跨平台性能可移植性上还有局限，不容易达到 Transformer 关键算子高性能自动代码生成的目的。针对上述问题本文提出 DeepGen --一个基于多级中间表示（Multi-Level Intermediate Representation, MLIR）的高性能算子生成框架，解决 Transformer 模型关键算子遇到的以上问题。本文的主要贡献是：

(1) 本文提出了一种线程块级与线程级相结合的双层中间表示（Intermediate Representation, IR）及其自动降级方法。为了解决现有编译器在代码生成过程中难以保留块级语义、无法显式表达计算和内存访问结构的问题，本文设计了基于 MLIR 的 DeepGen 方言（Dialect），对矩阵乘法、多头注意力机制中的线程块级计算逻辑、数据访存和并行映射关系进行统一建模。该方法通过显式定义 Alloc Buffer、Copy、Matmul 和 Block 等操作（Operation），实现计算逻辑与数据访存的解耦；同时结合算子计算方式和缓冲区（Buffer）依赖关系，自动推断线程块及线程束内部的索引布局，完成线程块级 IR 向低级线程级 IR 的降级转换。该贡献使算子语义能够在不同抽象层级间正确传递，为跨平台性能移植、内存访问优化、并行调度和底层代码生成提供了统一且清楚的表示基础。

(2) 本文提出了一套针对异构平台的跨平台优化 Pass 集合以及自动调优方法。优化 Pass 主要采用设计软件流水线（Software Pipeline）、局部规约轴切分（Local Split K）、线程块访存重排（Block Swizzle）和离散化访存（Access Scatter）等技术，来对不同硬件平台进行优化。优化 Pass 在计算调度、规约拆分、线程块映射、共享存储访问等多个方面进行了处理，进而提高所生成代码的执行效率。其次为解决大规模参数空间和手动调优成本高的问题，本文构建了一个依靠空间剪枝的自动调优框架。该框架依靠生成候选参数空间和剔除无效配置，自动进行面向目标硬件的优化执行配置搜索，从而使算子生成的自动化程度得到提高。

(3) 本文在 NVIDIA A100 和 V100 GPU 与海光 Z100 和 K100 DCU 的异构平台上，对 DeepGen 进行了实验验证。实验结果表明，在进行矩阵乘法测试时，DeepGen 比 cuBLAS 和 rocBLAS 分别达到了最高 1.26 倍和 1.24 倍的性能提升。在注意力机制测试里，相较 Triton 在 NVIDIA 和海光平台上的性能提高最高分别达到 4.12 倍和 4.87 倍。端到端推理实验结果说明，DeepGen 在模型 BERT、GPT-2 和 Llama-2 上都明显提高性能，具有不错的工程应用价值、跨平台适应能力。数值精度验证结果说明，DeepGen 生成的算子输出和 PyTorch 基准测试结果非常一致，表明本文提出的方法可以提高性能，并且能够保证数值正确性和工程可靠性。
}
%中文关键词%
\ckeywords{Transformer；MLIR；算子自动代码生成；异构计算；性能可移植性}

%英文摘要%
\eabstract{

The Transformer architecture has been widely used in natural language processing, code generation, and large model inference. Matrix multiplication (Matmul) and multi-head attention (MHA), as two key operators in Transformer models, have a direct impact on model deployment performance. However, heterogeneous hardware platforms differ significantly in execution models, memory hierarchies, and software ecosystems. Traditional approaches based on vendor-provided high-performance libraries or manually optimized kernels often suffer from complex adaptation processes, high development costs, and limited performance portability. In addition, existing deep learning compilers still have limitations in block-level semantic representation, low-level memory access control, and cross-platform performance portability, making it difficult to automatically generate high-performance code for key Transformer operators. To address these issues, this thesis proposes DeepGen, an MLIR-based framework for high-performance operator generation. The main contributions of this work are as follows.

First, this thesis proposes a two-level intermediate representation (IR) method and an automatic lowering method that combine thread-block-level and thread-level IR. To address the difficulty of preserving block-level semantics and explicitly representing computation and memory access structures during code generation, the DeepGen dialect is designed based on MLIR. This dialect provides unified modeling of thread-block-level computation, data movement, and parallel mapping for matrix multiplication and multi-head attention. Operations such as Alloc Buffer, Copy, Matmul, and Block are defined to decouple computational logic from data movement. In addition, according to operator semantics and buffer dependencies, the index layout inside thread blocks and warps is automatically inferred to support the lowering from thread-block-level IR to thread-level IR. This method ensures that operator semantics are correctly propagated across different abstraction levels, laying a foundation for cross-platform performance portability, memory access optimization, parallel scheduling, and low-level code generation. 

Second, this thesis proposes a set of cross-platform optimization passes and an automatic tuning method for heterogeneous platforms. The optimization passes employ techniques such as Software Pipeline, Local Split K, Block Swizzle, and Access Scatter to optimize generated code for different hardware platforms. These passes improve execution efficiency in terms of computation scheduling, reduction splitting, thread block mapping, and shared memory access. To reduce the cost of manual tuning in a large parameter space, an automatic tuning framework based on search space pruning is further constructed. By generating candidate configurations and eliminating invalid ones, the framework automatically searches for optimized execution configurations on target hardware, thereby improving the automation of operator generation. 

Third, this thesis evaluates DeepGen on heterogeneous platforms, including NVIDIA A100 and V100 GPUs as well as Hygon Z100 and K100 DCUs. Experimental results show that, in matrix multiplication tests, DeepGen achieves up to 1.26 times and 1.24 times performance improvement over cuBLAS and rocBLAS, respectively. In the evaluation of attention operators, DeepGen achieves up to 4.12 times and 4.87 times performance improvement over Triton on NVIDIA and Hygon platforms, respectively. End-to-end inference experiments on BERT, GPT-2, and Llama-2 further demonstrate that DeepGen can significantly improve model execution performance. Numerical accuracy verification shows that the outputs of operators generated by DeepGen are highly consistent with PyTorch baseline results, indicating that the proposed method improves performance while ensuring numerical correctness and engineering reliability.
}
%英文关键词%
\ekeywords{Transformer; MLIR; automatic operator code generation; heterogeneous computing; performance portability}

\makecover

\clearpage
